% !PRINTED 81
\weberpage{81}
%<<CONT
so ist die Beziehung zwischen den $\alpha$ und den $\beta$ durch die
Gleichung

\weq{1}{\omega_1^{\alpha_1\beta_1}\,\omega_2^{\alpha_2\beta_2}\,\dots\,\omega_\mu^{\alpha_\mu\beta_\mu}=1}
ausgedrückt.

Die Invarianten der Gruppe $\mathfrak{B}$, die nach \S.\,10 lauter Primzahlpotenzen sind, wollen wir mit
\[
i_1, i_2, \dots i_\nu \quad (\text{Invarianten von } \mathfrak{B})
\]
bezeichnen. Wir nehmen diese Invarianten als beliebig gegebene
Primzahlpotenzen an und bezeichnen mit $J$ ihr kleinstes gemeinschaftliches Vielfache, so dass
\weq{2}{J=i_1i_1'=i_2i_2'=\dots=i_\nu i_\nu'}
gesetzt werden kann.

Ausserdem soll $\mathfrak{A}$ als primärer Theiler von $\mathfrak{N}$ vorausgesetzt werden, was nach \S.\,19, 8. mit der Bedingung gleichbedeutend ist:

\gesperrt{IV. Keiner der Indices $\beta$ soll für alle Zahlen $b$, wenn die entsprechende Primzahl $q$ mehrmals in $m$ aufgeht, durch $q$, oder wenn $q$ nur einmal in $m$ aufgeht, durch $q-1$ theilbar sein.}

Nach der Definition der Invarianten muss die Gruppe $\mathfrak{B}$ eine Basis haben, deren Elemente
\[
\begin{aligned}
&\phantom{\text{von den Graden}} \quad g_1, g_2, \dots g_\nu \quad (\text{Basis von } \mathfrak{B})\\
&\text{von den Graden} \quad i_1, i_2, \dots i_\nu
\end{aligned}
\]
sind, so dass jede Zahl $b$ einmal und nur einmal in der Form
\weq{3}{b\equiv g_1^{x_1}g_2^{x_2}\dots g_\nu^{x_\nu}\ (\mathrm{mod}\ m)}
enthalten ist, wenn die Exponenten $x_1, x_2, \dots x_\nu$ je ein volles Restsystem nach den Moduln $i_1, i_2, \dots i_\nu$ durchlaufen.

Alle Grössen von der Form (3) bilden, wenn $g_1, g_2, \dots g_\nu$ beliebige Zahlen sind, bei unbeschränkter Veränderlichkeit der ganzzahligen Exponenten $x$ gewiss eine Gruppe. Sollen aber die $g_h$ die Elemente einer Basis dieser Gruppe, und ihre Grade $i_h$ sein, so darf $1\equiv g_1^{x_1}g_2^{x_2}\dots g_\nu^{x_\nu}\ (\mathrm{mod}\ m)$ nur dann erfüllt sein, wenn $x_h$ durch $i_h$ theilbar ist, also:

\gesperrt{1. Die nothwendige und hinreichende Bedingung dafür, dass $g_1, g_2, \dots g_\nu$ die Elemente einer Basis einer Gruppe $\mathfrak{B}$ von den Graden $i_1, i_2, \dots i_\nu$ seien, ist die, dass die Congruenz}
\weq{4}{g_1^{x_1}g_2^{x_2}\dots g_\nu^{x_\nu}\equiv 1\ (\mathrm{mod}\ m)}
%CONT>>
